\setcounter{section}{57} % tema número N-1
\section{Servicios}

Nombre completo: El procesamiento cooperativo y la arquitectura cliente-servidor. Arquitectura SOA.

\localtableofcontents

\subsection{Procesamiento cooperativo}

Contra los sistemas centralizados. El procesamiento distribuido utiliza componentes autónomos y heterogéneos que se ejecutan en unidades de hardware interconectadas. Cada elemento lógico trabaja en una tarea común.

Estos elementos se sincronizan a un reloj común. Tienen concurrencia global. Tienen tolerancia a fallos, de forma que un nodo fallando se ve transparente al resto de usuarios. Los sistemas son abiertos y heterogéneos. PErmitiendo no hacer vendor-locking.

Con esto se consigue transparencia, fiabilidad, rendimiento, escalabilidad y flexibilidad.

\subsection{Arquitectura cliente-servidor}

Es uno de los modelos mas extendidos. Se basa en la existencia de un sistema servidor encargado de proporcionar servicios a un cliente.

Normalmente se parte en tres componentes: presentación, lógica de negocio y base de datos. En algunos casos se queda separada la base de datos de presentación y negocio.

\subsection{Arquitectura SOA}

El objetivo de la arquitectura SOA (Service Oriented Architecture) es ofrecer una estrategia para integrar servicios, no necesariamente aplicaciones finales. OASIS define SOA como \textit{Un paradigma para organizar capacidades distribuidas y bajo el control de propietarios y dominios}.

Un servicio expone una funcionalidad definida y concreta, mientras que la aplicación orquesta la ejecución de servicios (los servicios pueden llamar a mas servicios). Un servicio puede enmascarar a un sistema Legacy preexistente. Facilitando su integración.

La arquitectura define servicios a utilizar, interacciones, tecnologías implementadas e interfaces.

Hay 8 principios base:

\begin{enumerate}
    \item Contrato de servicio
    \item Bajo acoplamiento
    \item Abstracción
    \item Reusabilidad
    \item Autonomía
    \item Sin estado
    \item Garantía de descubrimiento
    \item Preparados para usarse en composiciones
\end{enumerate}

Se aplica un sistema de capas para separar ciertos elementos. En la capa de sistemas operacionales hay sistemas légacy o no orientados a aplicaciones, de esta forma se pueden integrar. Los componentes de servicio aseguran los acuerdos a su nive y la calidad de servicios externos. Se soportan por servidores de aplicaciones con alta disponibilidad, balanceo, etc. La capa de servicios permite que estos sean expuestos e invocados por aplicaciones. En la capa de coreografía se establece el flujo para que los servicios interactuen conjuntamente hacia el exterior.

\subsubsection{Colaboración entre servicios}

\begin{description}
    \item[Orquestación] Un proceso controlado totalmente por una única entidad. Esta define completamente las interacciones entre servicios, y la lógica para conducir las interacciones. Toda la lógica del flujo de control está en un único punto.
    \item[Coreografía] Un proceso de coreografía no es controlado por un único participante y se comporta mas como un proceso público y no ejecutable.
\end{description}

\subsubsection{Soa + Web services}
Se puede contruir sobre servicios web como SOAP. Según el W3C se puede definir un servicio web como un sistema diseñado para soportar la interacción máquina-máquina a través de red. Normalmente cuenta con una interfaz descrita e en una interaz procesable por un sistema por la que se intercambian mensajes sobre http.

\subsubsection{Entreprise Service Bus}

Categoría de productos y middlewares basados en estándares de servicios que habilitan la creación de estructuras de servicios. Definen la infraestructura de integración y es donde reside el grueso de la comunicación SOA. Son a su vez servicios SOA, permitiéndole aprovecharse de esas ventajas.

El objeto de los ESB es también la integración de sistemas no servicios dentro de un SOA para su uso por parte de otras aplicaciones.